\begin{lemma}[Bogolyubov for compact abelian groups, see {\cite[Lemma 2.1]{ruzsa6}}]
\label{lem:bogolyubov_compact_group}
Let $G$ be a compact abelian group with Haar measure $\mu$ and let $A \subseteq G$ of positive measure. Then $A - A + A - A$ contains a Bohr-$(k, \eta)$ set where $k, \eta$ depends only on $\mu(A)$.
\end{lemma}

\subsection{Kernels}
A kernel on $G$ is a non-negative continuous function that satisfies $\int_G K \, d\mu =1$. In our case, we will utilize the kernels supported on given Bohr sets whose Fourier transforms are non-negative. For a kernel $K$, we write $\lVert \widehat{K} \rVert_1$ to denote $\sum_{\gamma \in \Gamma} | \widehat{K}(\gamma)|$.  

\begin{lemma}[cf. {\cite[Lemma 4.3]{bhmp}}]
\label{lem:kernel_a}
    Given a finite set $\Lambda \subset \Gamma$ and $\eta \in (0, 1/2]$, there exists a kernel $K$ satisfying the following:
    \begin{enumerate}
        \item $K \geq 0, \widehat{K} \geq 0$ and $\int_G K \, d \mu = \lVert K \rVert_1 = 1$,
        \item $\lVert \widehat{K} \rVert_1 = \lVert K \rVert_{\infty} \leq 1/(C_0\eta)^{|\Lambda|}$ for some absolute constant $0 < C_0 \leq 1$, and         
        \item $K$ vanishes outside the Bohr set $B( \Lambda; \eta)$.
    \end{enumerate}
Consequently, 
\begin{equation} \label{eq:mub}
\mu(B( \Lambda; \eta)) \geq (C_0\eta)^{|\Lambda|}.
\end{equation} 
\end{lemma}
We remark that the bound \eqref{eq:mub} can also be obtained from an elementary covering argument (see \cite{tao}). 
\begin{proof}
% \textcolor{red}{I added some details to this proof. Please check to make sure they are correct.}
    First, for each $\lambda \in \Lambda$, there exists a kernel $K_{\lambda}: G \to [0, \infty)$ satisfying the following properties:
\begin{enumerate}
\item $\lVert K_{\lambda} \rVert_1 = 1$,
\item $\widehat{K}_{\lambda} \geq 0$,
\item $K_{\lambda}$ is supported on $B(\{\lambda\}; \eta) = \{x \in G: |\lambda(x) - 1| < \eta\}$,
\item $\|K_\lambda\|_{\infty} = K_{\lambda}(0) \leq 1/(C_0 \eta)$ \text{ for some absolute constant $0 < C_0 \leq 1$}.
\end{enumerate}		
Indeed, let $B = B(\{\lambda\}; \frac{\eta}{2})$ and let $K_{\lambda} = \frac{1_B}{\mu(B)}*\frac{1_B}{\mu(B)}$. Clearly the first and second properties are satisfied. Additionally, $K_{\lambda}$ is supported on $B(\{\lambda\}; \frac{\eta}{2}) + B(\{\lambda\}; \frac{\eta}{2}) \subset B(\{\lambda\}; \eta)$. 

Concerning the last property, we have for every $x \in G$,
\[
    K_{\lambda}(x) = \sum_{\gamma \in \Gamma} \widehat{K}_{\lambda}(\gamma) \gamma(x)
\]
and so
\[
    \left| K_{\lambda}(x) \right| \leq \sum_{\gamma \in \Gamma} \widehat{K}_{\lambda}(\gamma) = K_{\lambda}(0).
\]
Therefore, $\lVert K_{\lambda} \rVert_{\infty} = K_{\lambda}(0) = \frac{1}{\mu(B)}$.
Since $\lambda$ is continuous, its image $\lambda(G)$ is a closed subgroup of $S^1 = \{ z \in \C: |z|=1\}$, and so it is either $S^1$ or $\{ z \in \C : z^q=1\}$ for some $q \in \N$. Since $\lambda$ is a homomorphism, it is measure-preserving (see Lemma \ref{lem:pushforward} below). Hence $\mu(B)$ is equal to the normalized Haar measure of the set
\[
\left\{ z \in S^1 : |z-1| < \frac{\eta}{2} \right\}
\]
in the group $\lambda(G)$. In either case, where $\lambda(G)=S^1$ or $\{ z \in \C : |z|^q=1\}$, we find that $\mu(B) \geq C_0 \eta$ for some absolute constant $0 < C_0 \leq 1$. Therefore, $\lVert K_{\lambda} \rVert_{\infty} \leq 1/(C_0 \eta)$.  
%\Hnote{Is $C=1$?} \Anote{No, it is slightly greater than $1$ since $\eta/2$ is the Euclidean distance while $\mu(B)$ corresponds to the distance on the unit circle.}
	
We now define
    \[
        \widetilde{K} = \prod_{\lambda \in \Lambda} K_{\lambda}.
    \]
    It follows that $\widetilde{K} \geq 0$ and $\widetilde{K}$ is supported on $B(\Lambda; \eta)$. Repeatedly using the fact that $\widehat{fg}(\gamma) = \sum_{\lambda \in \Gamma} \widehat{f}(\lambda) \widehat{g}(\gamma - \lambda)$ for all $f, g \in L^{\infty}(G)$, we have $\widehat{\widetilde{K}} \geq 0$. Likewise, since $\widehat{K}_{\lambda}(0) = \lVert K_{\lambda} \rVert_1 = 1$, we have $\lVert \widetilde{K} \rVert_1 = \int_G \widetilde{K} \ d \mu =  \widehat{\widetilde{K}}(0) \geq 1$.
    
    For every $x \in G$, $\widetilde{K}(x) = \sum_{\gamma \in \Gamma} \widehat{\widetilde{K}}(\gamma) \gamma(x)$ and so 
\[
    |\widetilde{K}(x)| \leq \sum_{\gamma \in \Gamma} |\widehat{\widetilde{K}}(\gamma)| = \lVert \widehat{\widetilde{K}} \rVert_1.
\]
It follows that $\lVert \widetilde{K} \rVert_{\infty} \leq \lVert \widehat{\widetilde{K}} \rVert_1$.
Moreover,
\[
    \widetilde{K}(0) = \sum_{\gamma \in \Gamma} \widehat{\widetilde{K}}(\gamma) \gamma(0) = \sum_{\gamma \in \Gamma} \widehat{\widetilde{K}}(\gamma) = \lVert \widehat{\widetilde{K}} \rVert_1
\]
because $\widehat{\widetilde{K}}(\gamma) \geq 0$ for all $\gamma$. Thus, $\lVert \widetilde{K} \rVert_{\infty} = \lVert \widehat{\widetilde{K}} \rVert_1$. Upon defining $K = \widetilde{K}/\lVert \widetilde{K} \rVert_{1}$,
    we obtain the desired kernel.
    % \textcolor{red}{We know $\widetilde{K}(0) \leq 1/(C_0 \eta)^{|\Lambda|}$. Why does $K(0)$ also satisfy this property? To this end, we need to show $\lVert \widetilde{K} \rVert_1 \geq 1$, but this might not be true. }
    % satisfies all desired properties except $\lVert \widehat{K} \rVert_1 = \lVert K \rVert_{\infty}$ which we now check. Since each $K_{\lambda}$ attains maximum at $0$, $K$ also attains maximum at $0$. Therefore,
    % \[
    %     \lVert K \rVert_{\infty} = K(0) = \sum_{\gamma \in \Gamma} \widehat{K}(\gamma) \gamma(0) = \sum_{\gamma \in \Gamma} \widehat{K}(\gamma) = \lVert \widehat{K} \rVert_1.
    % \]
\end{proof}


\subsection{Homomorphisms} We will often make use of the following facts about homomorphisms $G \rightarrow G$.

\begin{lemma}\label{lem:homo-solutions}
Let $\phi: G \rightarrow G$ be a continuous homomorphism such that $[G:\phi(G)] = m$ is finite. Then for any $\gamma \in \Gamma$, there are at most $m$ elements $\chi \in \Gamma$ such that $\gamma = \chi \circ \phi$.
\end{lemma}
\begin{proof}
It is easy to see that for each $\gamma \in \Gamma$, the set $S_\gamma := \{ \chi \in \Gamma : \gamma = \chi \circ \phi\}$ is either empty, or a coset of the group $S_0$. On the other hand, $S_0$ is the annihilator of the group $\phi(G)$, so by \cite[Theorem 2.1.2]{rudin}, it is isomorphic to $G/\phi(G)$, and hence has cardinality $m$.
\end{proof}

\begin{lemma} \label{lem:composition}
Let $\phi, \psi: G \rightarrow G$ be  homomorphisms such that $[G:\phi(G)] = m$ and $[G:\psi(G)] = \ell$ are finite. Then
$[G: \phi ( \psi (G))] \leq m\ell$ is finite.
\end{lemma}
\begin{proof}
We have $[G:\phi(\psi(G))] = [G:\phi(G)] [\phi(G): \phi(\psi(G))]$. It suffices to show that $[\phi(G): \phi(\psi(G))] \leq \ell$. 

Let $x_1 + \psi(G), \ldots, x_\ell + \psi(G)$ be all cosets of $\psi(G)$ in $G$. Then $\phi(x_1) + \phi(\psi(G)), \ldots, \phi(x_\ell) + \phi(\psi(G))$ are all cosets of $\phi(\psi(G))$ in $\phi(G)$ (these are not necessarily distinct, so the actual number of cosets may be less than $\ell$), proving the desired claim. 
\end{proof}

\begin{lemma}
\label{lem:pushforward}
    Let $G, H$ be compact abelian groups and $\mu, \nu$ be the normalized Haar measures of $G$ and $H$, respectively. Suppose $\phi: G \to H$ is a continuous surjective homomorphism. Then $\phi_* \mu = \nu$ (i.e. $\nu(B) = \mu( \phi^{-1}(B))$ for any Borel set $B \subset H$).
\end{lemma}
\begin{proof}
   %\textcolor{red}{to be filled}.\Anote{Done.}
    Let $\nu_0 = \phi_* \mu$. By the uniqueness of the normalized Haar measure, it suffices to show that $\nu_0$ is a translation-invariant probability measure on $H$. First, $\nu_0$ is a probability measure because $\nu_0(H) = \mu(\phi^{-1}(H)) = \mu(G) = 1$. Now let $B \subset H$ be a Borel set and $h_0 \in H$ be arbitrary. Since $\phi$ is surjective, there exists $g_0 \in G$ such that $\phi(g_0) = h_0$. For any $g \in \phi^{-1}(B + h_0)$, we have
    \[
        \phi(g - g_0) = \phi(g) - \phi(g_0) \in B + h_0 - h_0 = B.
    \]
    Therefore, $\phi^{-1}(B + h_0) \subseteq \phi^{-1}(B) + g_0.$ On the other hand, 
    \[
        \phi(\phi^{-1}(B) + g_0) \subseteq B + h_0
    \]
    and so $\phi^{-1}(B + h_0) = \phi^{-1}(B) + g_0$. Since $\mu$ is translation-invariant on $G$, it follows that
    \[
        \nu_0(B + h_0) = \mu(\phi^{-1}(B + h_0)) = \mu(\phi^{-1}(B) + g_0) = \mu(\phi^{-1}(B)) = \nu_0(B).
    \]
    Thus $\nu_0$ is translation-invariant on $H$ and so $\nu_0 = \nu$. 
\end{proof}


\begin{lemma} \label{lem:finite-index}
Let $\phi:G \rightarrow G$ be a continuous homomorphism such that $[G:\phi(G)]=m$ is finite. Then for any measurable set $A \subset G$, we have 
\begin{equation}
\mu( A ) \leq m \mu( \phi(A) )
\end{equation}
and
\begin{equation}
\mu( \phi^{-1}(A) ) \leq m \mu( A ).
\end{equation}
Consequently, if $f \in L^1(G)$ is nonnegative, then
\[
\int_G f(x) \, d\mu(x) \geq \frac{1}{m} \int_G f( \phi(x) ) \, d\mu(x).
\]
\end{lemma}

\begin{proof}
First, since $\phi$ is continuous and $G$ is compact, $\phi(G)$ is a compact subgroup of $G$. Since $G$ is Hausdorff, $\phi(G)$ is closed. In other words, $\phi(G)$ is a closed subgroup of $G$. 

%First we show that $nG$ is closed. Indeed, let $(g_i)$ be a sequence in $G$ and $g = \lim_{i \to \infty} n g_i$. Since $G$ is compact, there exists a convergent subsequence $(g_i')$ of $(g_i)$. Let $h = \lim_{i \to \infty} g_i'$. Then by continuity, $nh = g$. Therefore $g \in nG$. Hence $nG$ is closed. Since any closed subset of compact set is compact, we have $nG$ is a compact subgroup of $G$ (\textcolor{red}{check!}).

Observe that since $G$ is partitioned into $m$-many translates of $\phi(G)$, $\mu(\phi(G)) = 1/m$.
Define $\lambda(B) = m \mu(B)$ for any Borel set $B \subseteq \phi(G)$. Now $\lambda$ is a translation-invariant probability measure on $\phi(G)$, and so it is equal to the normalized Haar measure on $\phi(G)$. By \cref{lem:pushforward}, $\lambda = \phi_* \mu$. This means that for any Borel set $B \subset \phi(G)$, we have $\mu(\phi^{-1}(B)) = m \mu(B)$.

Let $A$ be any Borel set in $G$. Since $A \subset \phi^{-1}(\phi(A))$, we have $\mu(A) \leq \mu(\phi^{-1}(\phi(A)) = m \mu(\phi(A))$, and the first assertion is proved. Applying the first assertion to the set $\phi^{-1}(A)$, we get the second assertion. 

The third assertion follows from the second one, and the fact that $f$ can be approximated by functions of the form $\sum_{i=1}^n c_i 1_{A_i}$ for Borel sets $A_i$ and $c_i \geq 0$.
\end{proof}

The next lemmas deal with images and preimages of Bohr sets under homomorphisms.
 
\begin{lemma} \label{lem:bohr-homo1}
Let $B \subset G$ be a Bohr-$(k, \eta)$ set and $\phi: G \rightarrow G$ be a continuous homomorphism. Then $\phi^{-1}(B)$ is also a Bohr-$(k, \eta)$ set.
\end{lemma}

\begin{proof}
If $B = \{ x \in G: |\gamma_i(x) -1| < \eta \textup{ for }i=1,  \ldots,  k\}$ is a Bohr-$(k, \eta)$ set, then $\phi^{-1}(B) = \{ x \in G: |\gamma_i \circ \phi (x) -1| < \eta \textup{ for }i=1,  \ldots,  k\}$ is also a Bohr-$(k, \eta)$-set.
\end{proof}


% \begin{lemma}[Lower bound for measure of Bohr sets]
% \label{lem:bigness_bohr_set}
% If $G$ be a compact abelian group with Haar measure $\mu$ then every Bohr-$(k, \eta)$ set $B$ of $G$ satisfies $\mu(B) \geq \eta^k$.
% \end{lemma}
% \begin{proof}
% \Anote{Need a proof or reference.} \Hnote{This is already contained in Lemma 13}
% \end{proof}



The next lemma is more surprising.

\begin{lemma}[cf. Griesmer {\cite[Lemma 1.7] {griesmer-br}}] \label{lem:bohr-homo2}
Let $B \subset G$ be a Bohr-$(k, \eta)$ set and $\phi: G \rightarrow G$ be a continuous homomorphism such that $[G:\phi(G)] = m < \infty$. Then $\phi(B)$ contains a Bohr-$(k', \eta')$ set, where $k', \eta'$ depend on $k, \eta$ and $m$.
\end{lemma}

\begin{proof}
% \Hnote{Find proof that does not use Bogolyubov!!!!!!}
% \Anote{Here is a proof using Bogolyubov.}
Suppose $B = \{x \in G: |\gamma_i(x) - 1| < \eta \text{ for } 1 \leq i \leq k\}$ where $\gamma_i \in \Gamma$. Then 
\[
    A = \{x \in G: |\gamma_i(x) - 1| < \eta/4 \text{ for } 1 \leq i \leq k\}
\]
satisfies $A - A + A - A \subseteq B$. The bound \eqref{eq:mub} implies that $\mu(A) \geq (C_0\eta/4)^k$ for some absolute constant $C_0>0$. 

In view of \cref{lem:finite-index}, $\mu(\phi(A)) \geq \mu(A)/m \geq \frac{(C_0 \eta)^k}{4^k m}$. Therefore, by \cref{lem:bogolyubov_compact_group}, the set $\phi(B) \supseteq \phi(A) - \phi(A) + \phi(A) - \phi(A)$ is a Bohr-$(k', \eta')$ set where $k', \eta'$ depend only on $\mu(\phi(A))$, which is bounded below by $\frac{(C_0 \eta)^k}{4^k m}$.
\end{proof}

\subsection{Counting lemmas} 
% \textcolor{red}{Do we want to keep $\lVert K \rVert_{A(G)}$ or $\lVert K \rVert_1$?}

% For a function $K: G \to \R$, we define
% \begin{equation}
%     \lVert K \rVert_{A(G)} = \sum_{\gamma \in \Gamma} \left|\widehat{K}(\gamma)\right|
% \end{equation}
% and set $A(G) = \{K \in L^1(G): \lVert K \rVert_{A(G)} < \infty\}$.

\begin{lemma}[cf. {\cite[Lemma 2]{bourgain-roth}}]
\label{lem:bourgain_lemma_2}
 Let $\phi, \psi: G \rightarrow G$ be continuous homomorphisms such that $\phi(G), \psi(G)$ have finite index in $G$. Then for $f_1, f_2, f_3 \in L^\infty(G)$ and $K \in L^1(G)$ such that $\widehat{K} \in L^1(\Gamma)$, we have
\begin{equation} \label{eq:k}
\left| \iint_{G^2} f_1(x)f_2(x+ \phi(y))f_3(x+ \psi(y)) K(y) \, d\mu(x) d\mu(y) \right| \ll \| \widehat{f_1} \|_{\infty} \|f_2 \|_2 \| f_3\|_2 \| \widehat{K} \|_1  
\end{equation}
where the implied constant depends only on the indexes of $\phi(G)$ and $\psi(G)$ in $G$.
\end{lemma}

\begin{proof}
Since linear combinations of characters are dense in $L^1(G)$, without loss of generality, we can assume $f_1, f_2, f_3$ and $K$ are equal to their Fourier series. For $x \in G$, write $g(x) = \int_{G} f_2(x+ \phi(y) )f_3(x+ \psi(y) ) K(y) \, d\mu(y)$.

By Parseval's theorem,
\begin{equation*} %\label{eq:k2}
\left| \int_G f_1(x) g(x) \, d\mu(x) \right|= \left| \sum_{\gamma \in \Gamma} \widehat{f_1}(\gamma) \widehat{g} (\overline{\gamma}) \right| \leq \| \widehat{f_1} \|_{\infty} \cdot \| \widehat{g} \|_1.
\end{equation*} 
%For the time being, we assume that the functions $f_2, f_3, K$ are equal to the sums of their Fourier series. 
Thus
\begin{eqnarray*}
g(x) &=& \int_{G} f_2(x+ \phi(y)) f_3(x+ \psi(y)) K(y) \, d\mu(y) \nonumber \\
&=& \int_G \left( \sum_{\gamma_2, \gamma_3, \gamma_0 \in \Gamma} \widehat{f_2}(\gamma_2) \gamma_2(x+ \phi(y)) \widehat{f_3}(\gamma_3) \gamma_3(x+ \psi(y)) \widehat{K}(\gamma_0) \gamma_0(y) \right) \, d\mu(y) \nonumber \\
&=& \int_G \left( \sum_{\gamma_2, \gamma_3, \gamma_0 \in \Gamma} \widehat{f_2}(\gamma_2) \widehat{f_3}(\gamma_3) \widehat{K}(\gamma_0) (\gamma_2 + \gamma_3)(x) (\gamma_2 \circ \phi +  \gamma_3 \circ \psi + \gamma_0 ) (y)\right) \, d\mu(y) \nonumber \\
&=&  \sum_{\substack{ \gamma_2, \gamma_3, \gamma_0 \in \Gamma,\\ \gamma_2 \circ \phi + \gamma_3 \circ \psi + \gamma_0=0}} \widehat{f_2}(\gamma_2) \widehat{f_3}(\gamma_3) \widehat{K}(\gamma_0) (\gamma_2 + \gamma_3)(x) %\label{eq:k3}.
\end{eqnarray*}
Consequently,
\begin{equation*} %\label{eq:k4}
\widehat{g}(\gamma) = \sum_{\substack{ \gamma_2, \gamma_3, \gamma_0 \in \Gamma,\\ \gamma_2 \circ \phi + \gamma_3 \circ \psi + \gamma_0=0, \\ \gamma_2 + \gamma_3 = \gamma}} 
\widehat{f_2}(\gamma_2) \widehat{f_3}(\gamma_3) \widehat{K}(\gamma_0) 
\end{equation*}
and
\begin{equation*} %\label{eq:k5}
\lVert \widehat{g} \rVert_1 \leq \sum_{\substack{ \gamma_2, \gamma_3, \gamma_0 \in \Gamma,\\ \gamma_2 \circ \phi + \gamma_3 \circ \psi + \gamma_0=0}} |\widehat{f_2}(\gamma_2) | \cdot | \widehat{f_3}(\gamma_3) | \cdot 
| \widehat{K}(\gamma_0) |.
%&=& \sum_{\gamma_2, \gamma_3 \in \Gamma} |\widehat{f_2}(\gamma_2) | \cdot | \widehat{f_3}(\gamma_3) | \cdot 
%| \widehat{K}(-\gamma_2 - 2 \gamma_3) |
\end{equation*}
Therefore, it suffices to show that for each $\gamma_0 \in \Gamma$, we have
\[
\sum_{\substack{ \gamma_2, \gamma_3 \in \Gamma,\\ \gamma_2 \circ \phi + \gamma_3 \circ \psi + \gamma_0=0}} |\widehat{f_2}(\gamma_2) | \cdot | \widehat{f_3}(\gamma_3) | \ll \|f_2 \|_2 \cdot \| f_3 \|_2,
\]
where the implicit constant depends only on the indexes $[G:\phi(G)]$ and $[G:\psi(G)]$.
By the Cauchy-Schwarz inequality and the Parseval's theorem, the left hand side is at most
\begin{eqnarray}
& & \| f_2 \|_2 \cdot \left( \sum_{\gamma_2} \left( \sum_{\substack{\gamma_3 \\ \gamma_3 \circ \psi = - \gamma_0 - \gamma_2 \circ \phi}} | \widehat{f_3}(\gamma_3) | \right)^2  \right)^{1/2} \nonumber \\
& \ll & \| f_2 \|_2 \cdot \left( \sum_{\gamma_2} \sum_{\substack{\gamma_3 \\ \gamma_3 \circ \psi = - \gamma_0 - \gamma_2 \circ \phi}} | \widehat{f_3}(\gamma_3) |^2  \right)^{1/2} \label{eq:index1} \\
& \ll & \| f_2 \|_2 \cdot \left( \sum_{\gamma_3}  | \widehat{f_3}(\gamma_3) |^2  \right)^{1/2} \label{eq:index2} \\
&=& \| f_2 \|_2 \cdot \| f_3 \|_2. \nonumber
\end{eqnarray}
In \eqref{eq:index1}, we use the fact that for each $\xi \in \Gamma$, there are at most $[G: \psi(G)]$ values of $\gamma_3$ such that $\gamma_3 \circ \psi = \xi$. Likewise, in \eqref{eq:index2}, we use the fact that for each $\xi \in \Gamma$, there are at most $[G: \phi(G)]$ values of $\gamma_2$ such that $\gamma_2 \circ \phi = \xi$. Both of these facts follow from \cref{lem:homo-solutions}.
\end{proof}

